\section*{Bab \ref{ch:g9:trig} --- Trigonometri pada Segitiga Siku-siku}

\begin{solution}{exo:g9:trig:1}
Sisi miringnya: $\sqrt{9^2 + 12^2} = \sqrt{81 + 144} = \sqrt{225} = 15$.

Sisi siku-siku yang lain: $\sqrt{17^2 - 8^2} = \sqrt{289 - 64} = \sqrt{225} = 15$.
\end{solution}

\begin{solution}{exo:g9:trig:2}
$7^2 + 24^2 = 49 + 576 = 625 = 25^2$: ya, siku-siku (menurut konvers
Pythagoras), dengan \omterm{def:g3:shapes:rightangle}{sudut siku-sikunya} di seberang sisi $25$.

$5^2 + 6^2 = 61 \neq 64 = 8^2$: bukan siku-siku.
\end{solution}

\begin{solution}{exo:g9:trig:3}
Sisi miringnya adalah $[EF]$ (yang di seberang \omterm{def:g3:shapes:rightangle}{sudut siku-siku} di
$D$). Sisi depan $\theta = \widehat E$ adalah $[DF]$; sisi sampingnya
adalah $[DE]$. Karena itu
\[
\cos\theta = \frac{DE}{EF}, \qquad
\sin\theta = \frac{DF}{EF}, \qquad
\tan\theta = \frac{DF}{DE}.
\]
\end{solution}

\begin{solution}{exo:g9:trig:4}
Sisi depannya: $10 \sin 28^\circ \approx 10 \times 0.469 = 4.7$.
Sisi sampingnya: $10 \cos 28^\circ \approx 10 \times 0.883 = 8.8$.
(Periksa: $4.7^2 + 8.8^2 \approx 22 + 77 \approx 100$.)
\end{solution}

\begin{solution}{exo:g9:trig:5}
$\tan\theta = \dfrac{4.5}{6} = 0.75$, jadi
$\theta = \tan^{-1}(0.75) \approx 37^\circ$.
\end{solution}

\begin{solution}{exo:g9:trig:6}
Pada \omterm{def:g6:shapes:triangles}{segitiga siku-siku} yang dibentuk pengamatnya, titik di tanah
di bawah dronenya, dan dronenya: ketinggian $120$ m adalah sisi
depan sudut elevasinya, dan jarak mendatar $d$ adalah sisi sampingnya.
Jadi
\[
\tan 32^\circ = \frac{120}{d},
\qquad
d = \frac{120}{\tan 32^\circ} \approx \frac{120}{0.625} = 192 \text{ m}.
\]
\end{solution}

\begin{solution}{exo:g9:trig:7}
Kenaikannya ($1.2$ m) adalah sisi depan sudutnya; panjang jalur
landainya $L$ adalah sisi miringnya:
\[
\sin 5^\circ = \frac{1.2}{L},
\qquad
L = \frac{1.2}{\sin 5^\circ} \approx \frac{1.2}{0.0872} \approx 13.8
\text{ m}.
\]
Jalur landainya harus sepanjang sedikitnya $13.8$ m.
\end{solution}

\begin{solution}{exo:g9:trig:8}
Dari $\cos^2\theta + \sin^2\theta = 1$:
$\sin^2\theta = 1 - 0.36 = 0.64$, dan $\sin\theta > 0$ untuk sudut
lancip, jadi $\sin\theta = 0.8$. Lalu
$\tan\theta = \dfrac{\sin\theta}{\cos\theta} = \dfrac{0.8}{0.6} =
\dfrac43$.
\end{solution}

\begin{solution}{exo:g9:trig:9}
\emph{\omterm{def:g6:shapes:triangles}{Segitiga sama kaki} siku-siku} dengan sisi siku-siku $1$: sisi
miringnya $\sqrt{1 + 1} = \sqrt2$; tiap sudut lancipnya $45^\circ$, dan
\[
\cos 45^\circ = \frac{1}{\sqrt2} = \frac{\sqrt2}{2} = \sin 45^\circ,
\qquad
\tan 45^\circ = \frac11 = 1 .
\]

\emph{\omterm{def:g6:shapes:triangles}{Segitiga sama sisi}} bersisi $1$: tingginya membelahnya menjadi
dua \omterm{def:g6:shapes:triangles}{segitiga siku-siku} dengan sisi miring $1$, alas $\frac12$, dan
tinggi $h = \sqrt{1 - \frac14} = \frac{\sqrt3}{2}$ (menurut
Pythagoras). Sudutnya adalah $60^\circ$ (di alasnya) dan $30^\circ$
(di puncaknya). Dengan membaca perbandingannya:
\[
\cos 60^\circ = \frac{1/2}{1} = \frac12, \quad
\sin 60^\circ = \frac{\sqrt3/2}{1} = \frac{\sqrt3}{2}, \quad
\cos 30^\circ = \frac{\sqrt3}{2}, \quad
\sin 30^\circ = \frac12 .
\]
\end{solution}

\begin{solution}{pb:g9:trig:1}
\textbf{1.} Tinggi di atas ketinggian mata adalah sisi depan
$35^\circ$, dengan sisi samping $60$ m:
$h = 60 \tan 35^\circ \approx 60 \times 0.700 = 42$ m.

\textbf{2.} $\tan\theta = 0.12$ memberi
$\theta = \tan^{-1}(0.12) \approx 6.8^\circ$ --- curam untuk sebuah
jalan, namun sudutnya sederhana. Jalan bersudut $45^\circ$ punya
$\tan 45^\circ = 1$: yaitu kemiringan $100\,\%$. (``Seratus
persen'' jauh dari tegak --- salah paham kegemaran di kereta gantung
ski.)

\textbf{3.} Pada \omterm{def:g6:shapes:triangles}{segitiga siku-siku} dengan sudut lancip $\theta$ dan
$90^\circ - \theta$, sisi depan yang satu adalah sisi samping yang
lain dan sisi miringnya dipakai bersama. Jadi
$\sin(90^\circ - \theta) = \frac{\text{sisi depannya}}{h} =
\frac{\text{sisi samping }\theta}{h} = \cos\theta$, dan secara
\omterm{def:g7:central:def}{simetris} $\cos(90^\circ - \theta) = \sin\theta$: yaitu
\emph{ko}\omterm{def:g9:trig:ratios}{sinus} adalah \omterm{def:g9:trig:ratios}{sinus} \emph{ko}mplemennya.

\textbf{4.} Tingginya: $8 \sin 60^\circ = 8 \times
\frac{\sqrt3}{2} = 4\sqrt3 \approx 6.93$ m. Kakinya:
$8 \cos 60^\circ = 4$ m persis.

\textbf{5.} $0.574^2 + 0.819^2 = 0.329 + 0.671 = 1.000$ (sampai
\omterm{def:g6:decimals:rounding}{pembulatannya}); dan $\left(\frac12\right)^2 +
\left(\frac{\sqrt3}{2}\right)^2 = \frac14 + \frac34 = 1$. Pasangan
nilai mana pun yang gagal memenuhi kesamaannya mengkhianati bacaan
yang salah atau kalkulator dalam mode yang salah --- pendeteksi galat
yang gratis.

\textbf{6.} $\tan 40^\circ = \dfrac hx$ dan
$\tan 30^\circ = \dfrac{h}{x + 200}$.

\textbf{7.} $x = \dfrac{h}{\tan 40^\circ}$ dan
$x + 200 = \dfrac{h}{\tan 30^\circ}$; dengan mengurangkannya,
$200 = h\left(\frac{1}{\tan 30^\circ} -
\frac{1}{\tan 40^\circ}\right)$, sehingga muncul rumusnya.
Angkanya: $\frac{1}{0.577} \approx 1.732$, $\frac{1}{0.839} \approx
1.192$, dengan \omterm{ex:g1:subtraction:difference}{selisih} $0.540$:
$h \approx \frac{200}{0.540} \approx 370$ m.

\textbf{8.} $x = \frac{h}{\tan 40^\circ} \approx
\frac{370}{0.839} \approx 441$ m. Periksa dari $A$:
$\frac{h}{x + 200} \approx \frac{370}{641} \approx 0.577 =
\tan 30^\circ$: jadi cocok.

\textbf{9.} Cara satu stasiun memerlukan jarak mendatar ke titik
yang \emph{tepat di bawah puncaknya} --- yang tidak terukur di
seberang jurang atau di dalam sebuah gunung. Cara
dua stasiun menggantinya dengan jarak yang dapat kamu langkahi
sendiri: yaitu $200$ m antara kedua titik pandangmu sendiri.

\textbf{10.} $h = \dfrac{200}{\sqrt3 - \frac{1}{\sqrt3}}
= \dfrac{200}{\frac{2}{\sqrt3}} = 100\sqrt3 \approx 173$ m ---
yaitu lantai kedua Menara Eiffel, yang terukur dengan dua pembidikan
dan sekali berjalan kaki.

\textbf{11.} Garis pandangnya $d$, \omterm{def:g3:shapes:shapes}{jari-jarinya} $R$ (ke titik
serempetnya, \omterm{def:g6:lines:perp}{tegak lurus} terhadap garis pandangnya) dan $R + h$ (ke
matanya): Pythagoras memberi $d^2 + R^2 = (R + h)^2 = R^2 + 2Rh + h^2$,
jadi $d^2 = 2Rh + h^2 \approx 2Rh$ ($h$ sangat kecil terhadap $R$).
Mata pada $1.80$ m:
$d \approx \sqrt{2 \times 6.37 \times 10^6 \times 1.8}
\approx 4\,800$ m --- jadi kaki langit di laut hanya lima kilometer
jauhnya.

\textbf{12.} $d \approx \sqrt{2 \times 6.37 \times 10^6 \times
324} \approx 64$ km. Patokan para pelaut: $3.6\sqrt{1.8} \approx
4.8$ km dan $3.6\sqrt{324} = 64.8$ km --- keduanya cocok.

\textbf{13.} $\tan\theta = \frac{2}{60}$ memberi $\theta \approx
1.9^\circ$. Bulan yang $0.5^\circ$ itu kira-kira \emph{\omterm{def:g3:division:half}{seperempat}}
ibu jari: Bulan tampak raksasa padahal disembunyikan sebuah kuku
--- mata adalah busur derajat yang buruk.

\textbf{14.} Tangganya: $\tan^{-1}\frac{17}{29} \approx 30^\circ$
--- jadi tangga yang nyaman mendaki hampir persis pada
$30^\circ$ milik nilai istimewanya. Pada $60^\circ$ dengan pijakan
$29$ cm, tinggi anak tangganya akan $29\tan 60^\circ = 29\sqrt3 \approx
50$ cm: itu tangga naik, bukan tangga rumah.

\textbf{15.} $\tan 10^\circ \approx 0.18$;
$\tan 45^\circ = 1$; $\tan 80^\circ \approx 5.7$;
$\tan 89^\circ \approx 57.3$. Ketika $\theta \to 90^\circ$, sisi
sampingnya mengerut menjadi tiada sedangkan sisi
depannya tetap bertahan:
jadi \omterm{def:g3:division:remainder}{hasil baginya} tumbuh melampaui setiap batas. Pada $90^\circ$
persis tidak ada segitiganya dan tidak ada nilainya --- grafik \omterm{def:g9:trig:ratios}{tangennya}
mendaki asimtot tegak, yang pertama dari banyak asimtot di buku
Sekolah Menengah.

\end{solution}
